\section[Desenvolvimento do Projeto]{Desenvolvimento do Projeto}

\subsection{Repositório do Código Fonte do Projeto}
  https://tools.ages.pucrs.br/certifica-o-por-meio-de-blockchain/blockchain-wiki

  https://tools.ages.pucrs.br/certifica-o-por-meio-de-blockchain/blockchain-api

  https://tools.ages.pucrs.br/certifica-o-por-meio-de-blockchain/blockchain-ui

  https://tools.ages.pucrs.br/certifica-o-por-meio-de-blockchain/ethereum-contract

  https://tools.ages.pucrs.br/certifica-o-por-meio-de-blockchain/ethereum-bridge

  Nosso projeto foi dividido em cinco repositórios diferentes. O primeiro e padrão da Wiki, o segundo do backend, o terceiro ao frontend , o quarto a API do contrato da nossa aplicação na etherium e por ultimo a "ponte" que liga a blockchain a nossa aplicação backend.

\subsection{Banco de Dados Utilizado}
 https://tools.ages.pucrs.br/planline/wiki/-/wikis/banco\_dados

 Para a nossa aplicação usamos o banco de dados relacional Postgrees. Para ele criamos quatro entidades, Certificado, Empresa, Usuário e Permissão. Além dos relcionamentos básicos de conexão entre as entidades criamos um levando em consideração a geração de logs. Além do Postgrees usamos o firebase para a autenticação de usuários.
\sloppy
\subsection{Arquitetura Utilizada}
 https://tools.ages.pucrs.br/certifica-o-por-meio-de-blockchain/blockchain-wiki/-/wikis/arquitetura

 https://tools.ages.pucrs.br/certifica-o-por-meio-de-blockchain/blockchain-wiki/-/wikis/infraestrutura

 Para nossa arquitetura, organizamos o sistema em três camadas distintas: o backend, o frontend e a blockchain. A interação entre o frontend e a blockchain ocorria exclusivamente através do backend. Isso significa que todas as requisições e operações realizadas pelo frontend relacionadas à blockchain passavam primeiro pelo backend, que atuava como intermediário. Essa abordagem permitiu maior controle, segurança e facilita a integração entre as diferentes partes do sistema, garantindo que as operações na blockchain sejam realizadas de maneira eficiente e segura, sem expor diretamente o frontend à complexidade da infraestrutura blockchain.

\subsection{Protótipos das Telas Desenvolvidas}
\sloppy
  https://tools.ages.pucrs.br/certifica-o-por-meio-de-blockchain/blockchain-wiki/-/wikis/design/mockups

  Para a criação de mockups, usamos o figma. Mais detalhes sobre os mockups, utilize o link acima.

\subsection{Tecnologias Utilizadas}
  https://tools.ages.pucrs.br/certifica-o-por-meio-de-blockchain/blockchain-wiki/-/wikis/arquitetura
  
  Para o frontend usamos TypeScript e  o framework React. Para o backend usamos o Java 21 em conjunto com o  framework SpringBoot. Para a blockchain, além de usarmos a Etherium, usamos sua linguagem Solidity para a criação de contratos.
